\documentclass[11pt]{article}
\usepackage{amssymb}
\usepackage{amsthm}
\usepackage{enumitem}
\usepackage{amsmath, physics}
\usepackage{bm}
\usepackage{adjustbox}
\usepackage{mathrsfs}
\usepackage{graphicx}
% \usepackage{siunitx}
\usepackage{physunits}
\usepackage[mathscr]{euscript}

\title{\textbf{Solved selected problems of Classical Electrodynamics - Hans Ohanian}}
\author{Franco Zacco}
\date{}

\addtolength{\topmargin}{-3cm}
\addtolength{\textheight}{3cm}

\newcommand{\Nat}{\mathbb{N}}
\newcommand{\Z}{\mathbb{Z}}
\newcommand{\Q}{\mathbb{Q}}
\newcommand{\R}{\mathbb{R}}
\newcommand{\diam}{\text{diam}}
\newcommand{\cl}{\text{cl}}
\newcommand{\bdry}{\text{bdry}}
\newcommand{\inter}{\text{int}}
\newcommand{\hatx}{\bm{\hat{x}}}
\newcommand{\haty}{\bm{\hat{y}}}
\newcommand{\hatz}{\bm{\hat{z}}}
\newcommand{\hatr}{\bm{\hat{r}}}
\newcommand{\hatn}{\bm{\hat{n}}}
\newcommand{\hatrho}{\bm{\hat{\rho}}}
\newcommand{\hatphi}{\bm{\hat{\phi}}}
\newcommand{\hattheta}{\bm{\hat{\theta}}}
\newcommand{\vecx}{\bm{x}}
\newcommand{\varep}{\varepsilon}


\theoremstyle{definition}
\newtheorem*{solution*}{Solution}
\renewcommand*{\proofname}{\bf{Solution}}

\begin{document}
\maketitle
\thispagestyle{empty}

\section*{Chapter 7 - Vector Calculus in Spacetime}

\subsection*{Exercises}

\begin{proof}{\textbf{Exercise 1.}}\\
We see that
\begin{align*}
    {a^\alpha}_\mu \eta_{\alpha\beta} {a^\beta}_\nu
    &= \begin{pmatrix}
        \gamma & -V\gamma & 0 & 0 \\
        -V\gamma & \gamma & 0 & 0 \\
        0 & 0 & 1 & 0 \\
        0 & 0 & 0 & 1
    \end{pmatrix}
    \begin{pmatrix}
        1 & 0 & 0 & 0 \\
        0 & -1 & 0 & 0 \\
        0 & 0 & -1 & 0 \\
        0 & 0 & 0 & -1
    \end{pmatrix}
    \begin{pmatrix}
        \gamma & -V\gamma & 0 & 0 \\
        -V\gamma & \gamma & 0 & 0 \\
        0 & 0 & 1 & 0 \\
        0 & 0 & 0 & 1
    \end{pmatrix} \\
    &= \begin{pmatrix}
        \gamma & V\gamma & 0 & 0 \\
        -V\gamma & -\gamma & 0 & 0 \\
        0 & 0 & -1 & 0 \\
        0 & 0 & 0 & -1
    \end{pmatrix}
    \begin{pmatrix}
        \gamma & -V\gamma & 0 & 0 \\
        -V\gamma & \gamma & 0 & 0 \\
        0 & 0 & 1 & 0 \\
        0 & 0 & 0 & 1
    \end{pmatrix} \\
    &= \begin{pmatrix}
        \gamma^2 - V^2\gamma^2 & -V\gamma^2 + V\gamma^2 & 0 & 0 \\
        -V\gamma^2 + V\gamma^2 & V^2\gamma^2 - \gamma^2 & 0 & 0 \\
        0 & 0 & -1 & 0 \\
        0 & 0 & 0 & -1
    \end{pmatrix} \\
    &= \begin{pmatrix}
        \gamma^2(1 - V^2) & 0 & 0 & 0 \\
        0 & \gamma^2(V^2 - 1) & 0 & 0 \\
        0 & 0 & -1 & 0 \\
        0 & 0 & 0 & -1
    \end{pmatrix} \\
    &= \begin{pmatrix}
        1 & 0 & 0 & 0 \\
        0 & -1 & 0 & 0 \\
        0 & 0 & -1 & 0 \\
        0 & 0 & 0 & -1
    \end{pmatrix} \\
    &= \eta_{\mu\nu}
\end{align*}
Where we used that $1 - V^2 = 1/\gamma^2$.
\end{proof}

\cleardoublepage
\begin{proof}{\textbf{Exercise 2.}}\\
Equation (1.38) states that
\begin{align*}
    a^{km} a^{kl} = \delta^{ml}
\end{align*}
Let us define ${a^{\mu}}_\nu$ as 
\begin{align*}
    {a^{\mu}}_\nu
    &= \begin{pmatrix}
        1 & 0 \\
        0 & {a^k}_l
    \end{pmatrix}
\end{align*}
Where ${a^k}_l$ is a spatial rotation matrix.\\
Now, using ${a^\alpha}_\mu \eta_{\alpha\beta} {a^\beta}_\nu = \eta_{\mu\nu}$
when $\mu = m$ and $\nu = l$ i.e. we take the spatial indices of the matrix
${a^\mu}_\nu$ we see that
\begin{align*}
    {a^\alpha}_m \eta_{\alpha\beta} {a^\beta}_l 
    &= {a^0}_m \eta_{00} {a^0}_l
    + {a^1}_m \eta_{11} {a^1}_l
    + {a^2}_m \eta_{22} {a^2}_l
    + {a^3}_m \eta_{33} {a^3}_l \\
    &= {a^1}_m \eta_{11} {a^1}_l
    + {a^2}_m \eta_{22} {a^2}_l
    + {a^3}_m \eta_{33} {a^3}_l \\
    &= {a^k}_m \eta_{kj} {a^j}_l
\end{align*}
Where we used that ${a^0}_m$ is 0 when $m$ is a spatial coordinate.
Also, we changed $\alpha$ to $k$ and $\beta$ to $j$ so symbolize that they can
only take the spatial coordinate indices.\\
Therefore, since $\eta_{kj} = -\delta_{kj}$ and $\delta_{kj} {a^j}_l = {a^k}_l$
we have that
\begin{align*}
    {a^k}_m \eta_{kj} {a^j}_l &= \eta_{ml} \\
    -{a^k}_m \delta_{kj} {a^j}_l &= -\delta_{ml}\\
    {a^k}_m {a^k}_l &= \delta_{ml}
\end{align*}
\end{proof}

\cleardoublepage
\begin{proof}{\textbf{Exercise 3.}}\\
We note that
\begin{align*}
    \eta^{\alpha\mu}\eta_{\mu\nu}
    = \begin{pmatrix}
        1 & 0 & 0 & 0 \\    
        0 & -1 & 0 & 0 \\
        0 & 0 & -1 & 0 \\
        0 & 0 & 0 & -1
    \end{pmatrix}
    \begin{pmatrix}
        1 & 0 & 0 & 0 \\    
        0 & -1 & 0 & 0 \\
        0 & 0 & -1 & 0 \\
        0 & 0 & 0 & -1
    \end{pmatrix}
    = \begin{pmatrix}
        1 & 0 & 0 & 0 \\    
        0 & 1 & 0 & 0 \\
        0 & 0 & 1 & 0 \\
        0 & 0 & 0 & 1
    \end{pmatrix}
    = {\delta^{\alpha}}_{\nu}
\end{align*}
Then from equation (13) we know that
\begin{align*}
    x_\mu = \eta_{\mu\nu}x^\nu
\end{align*}
So multiplying both sides by $\eta^{\alpha\mu}$ we get that
\begin{align*}
    \eta^{\alpha\mu}x_\mu = \eta^{\alpha\mu}\eta_{\mu\nu}x^\nu
    = {\delta^\alpha}_\nu x^\nu = x^\alpha
\end{align*}
If we rename the free index $\alpha$ to $\nu$ we get the equation (15)
\begin{align*}
    x^\nu = \eta^{\nu\mu}x_\mu
\end{align*}
Therefore, equation (15) agrees with equation (13).
\end{proof}

\cleardoublepage
\begin{proof}{\textbf{Exercise 4.}}\\
Let
\begin{align*}
    {a^\nu}_\alpha
    = \begin{pmatrix}
        \gamma & -V\gamma & 0 & 0 \\
        -V\gamma & \gamma & 0 & 0 \\
        0 & 0 & 1 & 0 \\
        0 & 0 & 0 & 1 \\
    \end{pmatrix}
\end{align*}
Then
\begin{align*}
    {a_\mu}^\beta
    &= \eta_{\mu\nu} {a^\nu}_\alpha \eta^{\alpha\beta} \\
    &= \begin{pmatrix}
        1 & 0 & 0 & 0 \\
        0 & -1 & 0 & 0 \\
        0 & 0 & -1 & 0 \\
        0 & 0 & 0 & -1
    \end{pmatrix} 
    \begin{pmatrix}
        \gamma & -V\gamma & 0 & 0 \\
        -V\gamma & \gamma & 0 & 0 \\
        0 & 0 & 1 & 0 \\
        0 & 0 & 0 & 1 \\
    \end{pmatrix}
    \begin{pmatrix}
        1 & 0 & 0 & 0 \\
        0 & -1 & 0 & 0 \\
        0 & 0 & -1 & 0 \\
        0 & 0 & 0 & -1
    \end{pmatrix} \\
    &= \begin{pmatrix}
        \gamma & -V\gamma & 0 & 0 \\
        V\gamma & -\gamma & 0 & 0 \\
        0 & 0 & -1 & 0 \\
        0 & 0 & 0 & -1
    \end{pmatrix}
    \begin{pmatrix}
        1 & 0 & 0 & 0 \\
        0 & -1 & 0 & 0 \\
        0 & 0 & -1 & 0 \\
        0 & 0 & 0 & -1
    \end{pmatrix} \\
    &= \begin{pmatrix}
        \gamma & V\gamma & 0 & 0 \\
        V\gamma & \gamma & 0 & 0 \\
        0 & 0 & 1 & 0 \\
        0 & 0 & 0 & 1
    \end{pmatrix}
\end{align*}
\end{proof}

\cleardoublepage
\begin{proof}{\textbf{Exercise 5.}}\\
Let us compute ${a^\alpha}_\nu {a_\alpha}^\mu$ as follows
\begin{align*}
    {a^\alpha}_\nu {a_\alpha}^\mu
    &= \begin{pmatrix}
        \gamma & -V\gamma & 0 & 0 \\
        -V\gamma & \gamma & 0 & 0 \\
        0 & 0 & 1 & 0 \\
        0 & 0 & 0 & 1 \\
    \end{pmatrix}
    \begin{pmatrix}
        \gamma & V\gamma & 0 & 0 \\
        V\gamma & \gamma & 0 & 0 \\
        0 & 0 & 1 & 0 \\
        0 & 0 & 0 & 1 \\
    \end{pmatrix} \\
    &= \begin{pmatrix}
        \gamma^2 -V^2 \gamma^2 & V\gamma^2 - V\gamma^2 & 0 & 0\\
        -V\gamma^2 + V \gamma^2 & -V^2\gamma^2 + \gamma^2 & 0 & 0\\
        0 & 0 & 1 & 0 \\
        0 & 0 & 0 & 1
    \end{pmatrix} \\
    &= \begin{pmatrix}
        \gamma^2(1 - V^2) & 0 & 0 & 0\\
        0 & \gamma^2(1 - V^2) & 0 & 0\\
        0 & 0 & 1 & 0 \\
        0 & 0 & 0 & 1
    \end{pmatrix} \\
    &= \begin{pmatrix}
        1 & 0 & 0 & 0\\
        0 & 1 & 0 & 0\\
        0 & 0 & 1 & 0 \\
        0 & 0 & 0 & 1
    \end{pmatrix}
\end{align*}
Also, if we compute ${a_\alpha}^\mu{a^\alpha}_\nu$ we get that
\begin{align*}
    {a_\alpha}^\mu{a^\alpha}_\nu
    &= \begin{pmatrix}
        \gamma & V\gamma & 0 & 0 \\
        V\gamma & \gamma & 0 & 0 \\
        0 & 0 & 1 & 0 \\
        0 & 0 & 0 & 1 \\
    \end{pmatrix}
    \begin{pmatrix}
        \gamma & -V\gamma & 0 & 0 \\
        -V\gamma & \gamma & 0 & 0 \\
        0 & 0 & 1 & 0 \\
        0 & 0 & 0 & 1 \\
    \end{pmatrix} \\
    &= \begin{pmatrix}
        \gamma^2 -V^2 \gamma^2 & -V\gamma^2 + V\gamma^2 & 0 & 0\\
        V\gamma^2 - V \gamma^2 & -V^2\gamma^2 + \gamma^2 & 0 & 0\\
        0 & 0 & 1 & 0 \\
        0 & 0 & 0 & 1
    \end{pmatrix} \\
    &= \begin{pmatrix}
        1 & 0 & 0 & 0\\
        0 & 1 & 0 & 0\\
        0 & 0 & 1 & 0 \\
        0 & 0 & 0 & 1
    \end{pmatrix}
\end{align*}
So, in both cases we get the identity matrix, this implies that ${a^\alpha}_\nu$
is the inverse of ${a_\alpha}^\mu$ and viceversa.
\end{proof}

\cleardoublepage
\begin{proof}{\textbf{Exercise 6.}}\\
Equation (9) states that
\begin{align*}
    {a^\alpha}_\mu \eta_{\alpha\beta} {a^\beta}_\nu = \eta_{\nu\mu}
\end{align*}
Then multiplying the equation by $\eta^{\nu\sigma}$ we get that
\begin{align*}
    {a^\alpha}_\mu \eta_{\alpha\beta} {a^\beta}_\nu \eta^{\nu\sigma}
    = \eta_{\nu\mu} \eta^{\nu\sigma}
\end{align*}
We proved in Exercise 3 that
$\eta_{\nu\mu} \eta^{\nu\sigma} = {\delta^\sigma}_\mu$ and we know because of
equation (20) that
\begin{align*}
    \eta_{\alpha\beta} {a^\beta}_\nu \eta^{\nu\sigma}
    = {a_\alpha}^\sigma
\end{align*}
So we get that
\begin{align*}
    {a^\alpha}_\mu {a_\alpha}^\sigma = {\delta^\sigma}_\mu
\end{align*}
Or 
\begin{align*}
    {a_\alpha}^\mu {a^\alpha}_\nu = {\delta^\mu}_\nu
\end{align*}
Where we used that ${a^\alpha}_\mu$ and ${a_\alpha}^\sigma$ are inverses of
each other so they commute as we proved in Exercise 5 and we renamed the indexes
$\sigma \to \mu$ and $\mu \to \nu$.

\end{proof}

\cleardoublepage
\begin{proof}{\textbf{Exercise 7.}}\\
Let us consider the product $x^\mu x^\nu$. Since the position vectors
transform as
$$x'^\mu = {a^\mu}_\alpha x^\alpha$$
Then we can write that
\begin{align*}
    x'^\mu x'^\nu = {a^\mu}_\alpha x^\alpha {a^\nu}_\beta x^\beta
    = {a^\mu}_\alpha {a^\nu}_\beta x^\alpha x^\beta 
\end{align*}
So the product $x^\mu x^\nu$ transforms as a tensor, and hence it's a tensor.
\end{proof}

\begin{proof}{\textbf{Exercise 8.}}\\
We note that
\begin{align*}
    {a^\beta}_\nu = {a^\alpha}_\nu {\delta_\alpha}^\beta
\end{align*}
From equation (23) we have that
\begin{align*}
    {a_\beta}^\mu {a^\beta}_\nu = {\delta^\mu}_\nu
\end{align*}
So we can write that
\begin{align*}
    {a_\beta}^\mu {a^\alpha}_\nu {\delta_\alpha}^\beta = {\delta^\mu}_\nu
\end{align*}
Therefore ${\delta^\mu}_\nu$ transforms as a tensor and hence it's a tensor.
\end{proof}

\begin{proof}{\textbf{Exercise 9.}}\\
From equation (9) we know that
\begin{align*}
    \eta_{\alpha\beta} = {a^\mu}_\alpha \eta_{\mu\nu} {a^\nu}_\beta
\end{align*}
So multiplying the equation by ${a_\rho}^\alpha {a_\sigma}^\beta$ we get that
\begin{align*}
    {a_\rho}^\alpha {a_\sigma}^\beta\eta_{\alpha\beta}
    = {a_\rho}^\alpha {a^\mu}_\alpha \eta_{\mu\nu} {a^\nu}_\beta {a_\sigma}^\beta
    = \delta_\rho^\mu \eta_{\mu\nu} \delta^\nu_\sigma
\end{align*}
Therefore
\begin{align*}
    \eta_{\rho\sigma} = {a_\rho}^\alpha {a_\sigma}^\beta\eta_{\alpha\beta}
\end{align*}
So $\eta_{\mu\nu}$ transforms as a tensor.
\end{proof}

\cleardoublepage
\begin{proof}{\textbf{Exercise 10.}}\\
We know that
\begin{align*}
    A'^\mu = {a^\mu}_\alpha A^\alpha
    \quad\text{and that}
    \quad
    B'_\mu = {a_\mu}^\beta B_\beta
\end{align*}
Then
\begin{align*}
    A'^\mu B'_\mu
    = {a^\mu}_\alpha A^\alpha {a_\mu}^\beta B_\beta
    = {a^\mu}_\alpha {a_\mu}^\beta A^\alpha  B_\beta
    = \delta^\beta_\alpha A^\alpha  B_\beta
    = A^\beta  B_\beta
\end{align*}
\end{proof}

\begin{proof}{\textbf{Exercise 11.}}\\
Let $B^\mu$ be a vector and $A_{\alpha\mu}$ be a tensor then
\begin{align*}
    A'_{\alpha\mu} B'^\mu
    &= {a_\alpha}^\sigma {a_\mu}^\beta A_{\sigma \beta} {a^\mu}_\rho B^\rho \\
    &= {a_\alpha}^\sigma {a_\mu}^\beta {a^\mu}_\rho A_{\sigma \beta} B^\rho \\
    &= {a_\alpha}^\sigma \delta^\beta_\rho A_{\sigma \beta} B^\rho \\
    &= {a_\alpha}^\sigma A_{\sigma \beta} B^\beta
\end{align*}
Where we used the laws of transformation of $B^\mu$ and $A_{\alpha\mu}$, and
we see that $A_{\alpha\mu}B^\mu$ transforms as a vector.
\end{proof}

\begin{proof}{\textbf{Exercise 12.}}\\
Let $T^{\mu\nu}$ be a tensor, we see that
\begin{align*}
    {T'^\mu}_\nu = {a^\mu}_\alpha {a_\nu}^\beta {T^\alpha}_\beta
\end{align*}
If we let $\nu = \mu$ then we have that
\begin{align*}
    {T'^\mu}_\mu
    = {a^\mu}_\alpha {a_\mu}^\beta {T^\alpha}_\beta
    = \delta_\alpha^\beta {T^\alpha}_\beta
    = {T^\alpha}_\alpha
\end{align*}
We see that ${T^\mu}_\mu$ remains unchanged under Lorentz transformations so
it is a scalar.
\end{proof}

\cleardoublepage
\begin{proof}{\textbf{Exercise 13.}}\\
Let us compujte $u_\mu u^\mu$ as follows
\begin{align*}
    u_\mu u^\mu &= \eta_{\mu\nu} u^\nu u^\mu\\
    &= \frac{c^2}{1 -v^2/c^2} - \frac{(dx/dt)^2}{1 - v^2/c^2}
    - \frac{(dy/dt)^2}{1 - v^2/c^2}
    - \frac{(dz/dt)^2}{1 - v^2/c^2} \\
    &= \frac{c^2}{1 -v^2/c^2} - \frac{v^2}{1 - v^2/c^2} \\
    &= \frac{c^2 - v^2}{1 -v^2/c^2} \\
    &= \frac{c^2(1 - v^2/c^2)}{1 -v^2/c^2} \\
    &= c^2
\end{align*}
Where we used that $v^2 = (dx/dt)^2 + (dy/dt)^2 + (dz/dt)^2$.
\end{proof}

\begin{proof}{\textbf{Exercise 14.}}\\
We know that $u_\mu u^\mu = c^2$ and that $p^\mu = m u^\mu$ so
\begin{align*}
    p_\mu p^\mu = m^2 u_\mu u^\mu = m^2 c^2
\end{align*}
Therefore
\begin{align*}
    p_\mu p^\mu = \frac{E^2}{c^2} - p_x^2 - p_y^2 - p_z^2 = m^2c^2
\end{align*}
\end{proof}

\cleardoublepage
\begin{proof}{\textbf{Exercise 15.}}\\
We know from the Lorentz transformations that
\begin{align*}
    E' &= \frac{E - VP_x}{\sqrt{1 - V^2/c^2}} &
    P_x' &= \frac{P_x - VE/c^2}{\sqrt{1 - V^2/c^2}}\\
    P_y' &= P_y & P_z' &= P_z
\end{align*}
Where the primed quantities are the energy and momentum measured in a frame
$S'$ which is moving with velocity $V$ with respect to an inertial frame $S$.
In our case, we get that $E_0 = E'$ so
\begin{align*}
    E - VP_x &= E_0 \sqrt{1 - V^2/c^2}
\end{align*}
Also, in our case, in $S'$ we have that $P_x' = P_y' = P_z' = 0$ so we get that
\begin{align*}
    P_x &= VE/c^2 & P_y = P_z = 0
\end{align*}
Therefore
\begin{align*}
    E(1 - V^2/c^2) &= E_0 \sqrt{1 - V^2/c^2} \\
    E &= E_0 \frac{\sqrt{1 - V^2/c^2}}{1 - V^2/c^2} \\
    E &= \frac{E_0}{\sqrt{1 - V^2/c^2}}
\end{align*}
Replacing $E$ we get that
\begin{align*}
    P_x &= \frac{E_0 V/c^2}{\sqrt{1 - V^2/c^2}}
\end{align*}
Hence
\begin{align*}
    \bm{P} &= \frac{E_0 \bm{V}/c^2}{\sqrt{1 - V^2/c^2}}
\end{align*}
\end{proof}

\cleardoublepage
\begin{proof}{\textbf{Exercise 16.}}\\
Let $\Phi(t,x) = f(ct + x)$ and let us call $v = ct + x$ then, we get that
\begin{align*}
    \pdv{\Phi}{t} &= \pdv{f(v)}{t} = \pdv{f(v)}{v}\pdv{v}{t} = c\pdv{f(v)}{v} \\
    \pdv{\Phi}{x} &= \pdv{f(v)}{x} = \pdv{f(v)}{v}\pdv{v}{x} = \pdv{f(v)}{v}
\end{align*}
Differentiating again we have that
\begin{align*}
    \pdv[2]{\Phi}{t} &= \pdv{t}(c\pdv{f(v)}{v})
    = c\pdv{v}(\pdv{f(v)}{t})
    = c^2\pdv[2]{f(v)}{v} \\
    \pdv[2]{\Phi}{x} &= \pdv{x}(\pdv{f(v)}{v})
    = \pdv{v}(\pdv{f(v)}{x})
    = \pdv[2]{f(v)}{v}
\end{align*}
Then
\begin{align*}
    \frac{1}{c^2}\pdv[2]{\Phi}{t} - \pdv[2]{\Phi}{x} &=
    \frac{1}{c^2}\bigg(c^2\pdv[2]{f(v)}{v}\bigg) - \pdv[2]{f(v)}{v}
    = 0
\end{align*}
Therefore $f(ct + x)$ is a solution to equation (96). If we try $f(ct - x)$
with $v = ct - x$ we see that
\begin{align*}
    \pdv{\Phi}{x} &= \pdv{f(v)}{x} = \pdv{f(v)}{v}\pdv{v}{x} = -\pdv{f(v)}{v}
\end{align*}
and
\begin{align*}
    \pdv[2]{\Phi}{x} &= \pdv{x}(-\pdv{f(v)}{v})
    = \pdv{v}(-\pdv{f(v)}{x})
    = \pdv{v}(-\bigg(-\pdv{f(v)}{v}\bigg))
    = \pdv[2]{f(v)}{v}
\end{align*}
This is the same result we get for $f(ct + x)$, therefore equation (96)
is also satisfied by $f(ct - x)$.

\end{proof}

\cleardoublepage
\begin{proof}{\textbf{Exercise 17.}}\\
From the Lorentz-transformation equations for the current four-vector, we know
that
\begin{align*}
    \rho' &= \frac{\rho - Vj_x /c^2}{\sqrt{1 - V^2/c^2}} \qquad
    j_x' = \frac{j_x - V\rho}{\sqrt{1 - V^2/c^2}}
\end{align*}
If a wire along the $x$ axis is electrically neutral but carries a current in
the $xyz$ frame then $\rho = 0$ and $j_x \neq 0$ then
\begin{align*}
    \rho' &= -\frac{Vj_x /c^2}{\sqrt{1 - V^2/c^2}} \qquad
    j_x' = \frac{j_x}{\sqrt{1 - V^2/c^2}}
\end{align*}
Since $\rho' \neq 0$ this implies that the wire will not be neutral in the
$x'y'z'$ frame.
\end{proof}
\end{document}